\chapter{生殖系统}\label{ch:g8:reproductive-systems}

\Cref{ch:g8:sexual-reproduction}需要\omterm{def:g8:sexual-reproduction:gametes}{配子}；
\cref{ch:g8:puberty}启动了它们的工厂。本
章介绍工厂本身：\omterm{def:g4:animal-reproduction:sexual}{雄性}与\omterm{def:g4:animal-reproduction:sexual}{雌性}
生殖系统——器官，各自产生的
\omterm{def:g8:sexual-reproduction:gametes}{配子}，以及\omterm{def:g4:animal-reproduction:sexual}{雌性}系统非凡的每月
节律。朴素的解剖，朴素地命名：这些是像
\omterm{def:g4:heart-and-blood:heart}{心脏}和肺那样的器官，本章正是讲它们的。

\section{雄性生殖系统}

\begin{definition}[雄性生殖系统]\label{def:g8:reproductive-systems:male}
\omterm{def:g4:animal-reproduction:sexual}{雄性}系统的器官：
\begin{itemize}
  \item 两个\emph{睾丸}\index{睾丸}，悬于体温之外，
        位于\emph{阴囊}中——是
        \omterm{def:g8:sexual-reproduction:gametes}{配子}工厂：自\omterm{def:g8:puberty:puberty}{青春期}起，它们源源不断地产出
        \emph{精子}\index{精子}，持续不断地
        数以亿计；
  \item 管道储存并输送精子，途中
        有腺体汇入，其液体滋养并输送
        精子——共同形成\emph{精液}；
  \item \emph{阴茎}\index{阴茎}，精液由此
        排出体外——同一器官的另一职能是
        输送\omterm{def:g7:removing-wastes:kidneys}{尿液}
        （见\cref{def:g7:removing-wastes:kidneys}）；一套
        \omterm{def:g7:heart-and-circulation:rooms}{瓣膜}结构使两种职能严格分开。
\end{itemize}
\end{definition}

\begin{example}[近看精子细胞]\label{ex:g8:reproductive-systems:sperm}
在\omterm{def:g5:first-look-at-cells:microscope}{显微镜}下
（见\cref{def:g5:first-look-at-cells:microscope}），\omterm{def:g8:reproductive-systems:male}{精子}
\omterm{def:g5:first-look-at-cells:cell}{细胞}是为远行而精简的\omterm{def:g5:first-look-at-cells:cell}{细胞}：头部
几乎全是\omterm{def:g6:cell-unit-of-life:parts}{细胞核}——父方的贡献，高度压缩
（见\cref{rem:g5:first-look-at-cells:nucleus}）——以及一条
摆动的尾巴。人体最小的\omterm{def:g5:first-look-at-cells:cell}{细胞}，为游动而构造，
备有数日所需的养分。生产自\omterm{def:g8:puberty:puberty}{青春期}起，
一直持续到\omterm{def:g3:stages-of-life:stages}{老年期}：\omterm{def:g4:animal-reproduction:sexual}{雄性}工厂是持续运转的。
\end{example}

\section{雌性生殖系统}

\begin{definition}[雌性生殖系统]\label{def:g8:reproductive-systems:female}
\omterm{def:g4:animal-reproduction:sexual}{雌性}系统的器官，均位于下腹部：
\begin{itemize}
  \item 两个\emph{卵巢}\index{卵巢}——\omterm{def:g8:sexual-reproduction:gametes}{配子}
        的储存库与工厂：与\omterm{def:g8:reproductive-systems:male}{睾丸}不同，它们自出生起
        就储备了未来的\emph{\omterm{def:g4:animal-reproduction:sexual}{卵细胞}}，
        并自\omterm{def:g8:puberty:puberty}{青春期}起每次释放一个；
  \item 两条\emph{输卵管}\index{输卵管}，
        有纤毛的管道，接住释放的\omterm{def:g4:animal-reproduction:sexual}{卵细胞}——这里是
        通常的相遇之处，
        \cref{prop:g8:sexual-reproduction:core}的第 2 步；
  \item \emph{\omterm{prop:g5:human-reproduction:plan}{子宫}}（即
        \cref{prop:g5:human-reproduction:plan}中所说的\omterm{prop:g5:human-reproduction:plan}{子宫}），其
        厚实的\omterm{def:g3:bones-and-muscles:muscle}{肌肉}壁和可再生的内膜，是
        \omterm{prop:g8:sexual-reproduction:core}{受精卵}的着床之处；
  \item \emph{阴道}，连接\omterm{prop:g5:human-reproduction:plan}{子宫}与
        外界——精液在此被接纳，也是
        \cref{ex:g5:human-reproduction:birth}中的产道。
\end{itemize}
\end{definition}

\begin{example}[近看卵细胞]\label{ex:g8:reproductive-systems:egg}
\omterm{def:g4:animal-reproduction:sexual}{卵细胞}在作用之外，各方面都与\omterm{def:g8:reproductive-systems:male}{精子}相反：
人体\emph{最大}的\omterm{def:g5:first-look-at-cells:cell}{细胞}——约十分之一
毫米，肉眼勉强可见——圆形，静止不动，
储备着足够旅程最初几日用的养分
（\cref{rem:g2:animals-grow-up:egg}中的卵黄原理，
在\omterm{def:g5:first-look-at-cells:cell}{细胞}尺度上）。一座工厂持续产出数以百万计；
另一座每月仅释放一个备足养分的\omterm{def:g5:first-look-at-cells:cell}{细胞}：
\cref{prop:g4:animal-reproduction:strategies}中的
数量与照料的取舍，写入\omterm{def:g8:sexual-reproduction:gametes}{配子}
本身。
\end{example}

\section{循环}

\begin{proposition}[月经周期]\label{prop:g8:reproductive-systems:cycle}
自\omterm{def:g8:puberty:puberty}{青春期}起至约五十岁，\omterm{def:g4:animal-reproduction:sexual}{雌性}系统运行
一套约\emph{28 天}的重复程序——即
\emph{月经周期}\index{月经周期}，受
\cref{def:g8:puberty:hormones}所述\omterm{def:g8:puberty:hormones}{激素}链调控：
\begin{enumerate}
  \item \omterm{prop:g5:human-reproduction:plan}{子宫}内膜重建，增厚且
        血量丰富——着床之所已备妥；
  \item 周期中段，约第 14 天，一侧\omterm{def:g8:reproductive-systems:female}{卵巢}
        释放一个\omterm{def:g4:animal-reproduction:sexual}{卵细胞}——\emph{排卵}\index{排卵}；
        \omterm{def:g8:reproductive-systems:female}{输卵管}接住它，约一天内可
        \omterm{def:g5:flower-to-fruit:steps}{受精}；
  \item 未\omterm{def:g5:flower-to-fruit:steps}{受精}：内膜不再需要，它
        分解，经阴道在数日内
        排出，即\emph{月经
        期}\index{月经期}——程序
        重新开始；
  \item \omterm{def:g5:flower-to-fruit:steps}{受精}：\omterm{prop:g8:sexual-reproduction:core}{受精卵}落入
        已准备好的内膜
        （见\cref{def:g5:human-reproduction:pregnancy}），
        循环暂停，\omterm{def:g5:human-reproduction:pregnancy}{怀孕}开始。
\end{enumerate}
\end{proposition}
\admitted

\begin{omfigure}
\begin{tikzpicture}[scale=0.94]
  % cycle wheel
  \draw[very thick, omDef!60] (0,0) circle (2.3);
  % day marks
  \foreach \a/\d in {90/1, 0/8, -90/14, 180/22} {
    \fill[omDef] (\a:2.3) circle (0.07);
  }
  \node[font=\footnotesize, above] at (90:2.45) {第 1 天};
  \node[font=\footnotesize, right] at (0:2.45) {第 8 天};
  \node[font=\footnotesize, below] at (-90:2.45) {第 14 天};
  \node[font=\footnotesize, left] at (180:2.45) {第 22 天};
  % phases
  \draw[very thick, omThm] (90:2.3) arc[start angle=90, end angle=35,
    radius=2.3];
  \node[font=\footnotesize, omThm, align=center] at (58:3.5)
    {月经期：\\ 内膜脱落};
  \node[font=\footnotesize, omMeth, align=center] at (-35:4.0)
    {内膜重建，\\ 增厚备用};
  \fill[omProp] (-90:2.3) circle (0.12);
  \node[font=\footnotesize, omProp, align=center] at (-90:3.55)
    {排卵：\\ 释放一个卵细胞};
  % center
  \node[font=\small, align=center] at (0,0.25)
    {循环：\\ 约 28 天};
  \node[font=\footnotesize, align=center] at (0,-0.75)
    {（受精使循环暂停：\\ 怀孕）};
  % arrow direction
  \draw[->, thick, omDef] (150:2.3) arc[start angle=150,
    end angle=120, radius=2.3];
\end{tikzpicture}

{\small 把\omterm{prop:g8:reproductive-systems:cycle}{月经周期}画成一个轮盘：经期打开它，
\omterm{prop:g5:human-reproduction:plan}{子宫}内膜重建，\omterm{prop:g8:reproductive-systems:cycle}{排卵}在中途冠以高潮——而
\omterm{def:g5:flower-to-fruit:steps}{受精}是唯一能让轮盘停下来的事件。}
\end{omfigure}

\begin{remark}[生活中的周期]\label{rem:g8:reproductive-systems:living}
实话实说：前几年周期不规则，
机体仍在调试；周期长短因人而异，
月与月之间也不同；经期可能伴有痉挛
（\omterm{prop:g5:human-reproduction:plan}{子宫}是一块\omterm{def:g3:bones-and-muscles:muscle}{肌肉}，
见\cref{def:g3:bones-and-muscles:muscle}，并且在脱落期间仍在活动），
温暖、运动，以及必要时校医的建议，都有帮助；
\omterm{def:g1:healthy-habits:hygiene}{卫生}巾、\omterm{def:g1:healthy-habits:hygiene}{卫生}棉条和月经杯只是经期的日常用品，
和药店里任何一排货架上的东西一样寻常。
这些都不是疾病：
周期不过是系统按其构造正常运行。
\end{remark}

\begin{method}[按功能并列比较两个系统]\label{met:g8:reproductive-systems:sides}
为免混淆，按功能对照：
\begin{enumerate}
  \item \omterm{def:g8:sexual-reproduction:gametes}{配子}工厂：\omterm{def:g8:reproductive-systems:male}{睾丸}——\omterm{def:g8:reproductive-systems:female}{卵巢}；
  \item \omterm{def:g8:sexual-reproduction:gametes}{配子}：\omterm{def:g8:reproductive-systems:male}{精子}持续产生，数以百万计——
        \omterm{def:g4:animal-reproduction:sexual}{卵细胞}每月仅此一个，并已备好营养；
  \item 会合场所：\omterm{def:g8:reproductive-systems:female}{输卵管}，\omterm{def:g8:reproductive-systems:male}{精子}从阴道出发，
        穿过\omterm{prop:g5:human-reproduction:plan}{子宫}游到此处；
  \item 住所：男性一侧没有——\omterm{prop:g5:human-reproduction:plan}{子宫}，
        每月准备一次，位于女性一侧；
  \item 时间表：男性自\omterm{def:g8:puberty:puberty}{青春期}起持续产生，女性周期性，
        自\omterm{def:g8:puberty:puberty}{青春期}至约五十岁。
\end{enumerate}
\end{method}

\section{练习}

\begin{exercise}[$\star$]\label{exo:g8:reproductive-systems:1}
画出男性系统：说出每个器官的名称及其功能。
\end{exercise}

\begin{exercise}[$\star$]\label{exo:g8:reproductive-systems:2}
同样画出女性系统。
\end{exercise}

\begin{exercise}[$\star$]\label{exo:g8:reproductive-systems:3}
对比\omterm{def:g8:reproductive-systems:male}{精子}与\omterm{def:g4:animal-reproduction:sexual}{卵细胞}：大小、结构、数量，
以及营养储备。
\end{exercise}

\begin{exercise}[$\star$]\label{exo:g8:reproductive-systems:4}
\omterm{prop:g8:reproductive-systems:cycle}{排卵}时发生了什么，大约在哪一天，
\omterm{def:g4:animal-reproduction:sexual}{卵细胞}可\omterm{def:g5:flower-to-fruit:steps}{受精}的时间有多长？
\end{exercise}

\begin{exercise}[$\star$]\label{exo:g8:reproductive-systems:5}
从\omterm{prop:g5:human-reproduction:plan}{子宫}内膜的角度看，什么是\omterm{prop:g8:reproductive-systems:cycle}{月经期}——
它的到来说明了该周期的什么情况？
\end{exercise}

\begin{exercise}[$\star$]\label{exo:g8:reproductive-systems:6}
哪一事件使周期暂停，随后
\omterm{prop:g5:human-reproduction:plan}{子宫}中发生什么？
\end{exercise}

\begin{exercise}[$\star\star$]\label{exo:g8:reproductive-systems:7}
\cref{prop:g8:sexual-reproduction:core}中的三个
步骤在本章图中分别发生在哪些器官上？
\end{exercise}

\begin{exercise}[$\star\star$]\label{exo:g8:reproductive-systems:8}
把\omterm{def:g8:sexual-reproduction:gametes}{配子}视作
\cref{prop:g4:animal-reproduction:strategies}中两种策略的
缩影。哪一方承载哪种
逻辑？
\end{exercise}

\begin{exercise}[$\star\star$]\label{exo:g8:reproductive-systems:9}
为什么周期的前几年常常不规则，
而长度不一为何也正常？（两条依据见
\cref{rem:g8:reproductive-systems:living}。）
\end{exercise}

\begin{exercise}[$\star\star$]\label{exo:g8:reproductive-systems:10}
解释经期痉挛，
结合\cref{def:g3:bones-and-muscles:muscle}——并说出
两种有助缓解的方法。
\end{exercise}

\begin{exercise}[$\star\star$]\label{exo:g8:reproductive-systems:11}
两个工厂的时间表截然相反——
一方持续产生，另一方周期性产生并自出生起即有储备。把
每一方与其\omterm{def:g8:sexual-reproduction:gametes}{配子}的结构和作用联系起来。
\end{exercise}

\begin{exercise}[$\star\star\star$]\label{exo:g8:reproductive-systems:12}
“周期是\omterm{prop:g5:human-reproduction:plan}{子宫}每月提出的问题，
只有\omterm{def:g5:flower-to-fruit:steps}{受精}才能回答。”
把这一比喻展开为
\cref{prop:g8:reproductive-systems:cycle}中的四个编号事件。
\end{exercise}

\section{问题：诊所的教学模型}

\begin{problem}[{周末问题——用两个模型和一本日历
讲解两个系统}]\label{pb:g8:reproductive-systems:1}
一家健康中心的开放日使用三件教学道具：
男性系统模型、女性系统模型，以及一本大型
28 天周期日历。你是志愿讲解员。

\medskip\noindent\textbf{第一部分——两个模型。}
\begin{enumerate}
  \item 参观者沿男性模型，从\omterm{def:g8:reproductive-systems:male}{睾丸}走到
        出口。讲述路线，说明途中加入了哪些
        物质，以及一身两用器官的排列方式。
  \item 在女性模型上，指出\omterm{def:g8:sexual-reproduction:gametes}{配子}储备
        在哪里等待，\omterm{def:g5:flower-to-fruit:steps}{受精}通常发生在哪里，
        以及妊娠在哪里着床。
  \item 参观者问，两个模型中的\omterm{def:g8:sexual-reproduction:gametes}{配子}工厂
        为何差别如此之大——一方每天数以百万计，
        另一方每月仅一个。用
        \cref{ex:g8:reproductive-systems:egg}的结束句
        回答。
  \item 另一位参观者问，\omterm{def:g8:puberty:puberty}{青春期}之前
        两个系统在做什么。用一句话说明\cref{ch:g8:puberty}
        中的机制。
\end{enumerate}

\medskip\noindent\textbf{第二部分——日历。}
\begin{enumerate}[resume]
  \item 沿日历走一遍，标出经期、
        \omterm{prop:g5:human-reproduction:plan}{子宫}内膜重建、\omterm{prop:g8:reproductive-systems:cycle}{排卵}与易受孕日，
        并注明各自的大致日期。
  \item 有参观者以为第 28 天对每个人都一样。
        请纠正这一看法：
        见\cref{rem:g8:reproductive-systems:living}——
        两种变异。
  \item 在日历上标出，若\omterm{def:g5:flower-to-fruit:steps}{受精}发生在第 14 天左右——
        哪些事件会被取消，
        哪些事件会开始
        （\cref{def:g5:human-reproduction:pregnancy}
        延续这一过程的故事）。
  \item 一个青少年小声问，痉挛是否说明出了问题。
        给出\omterm{def:g3:bones-and-muscles:muscle}{肌肉}方面的答案和两条帮助办法，
        并说明超出这些办法时
        应去何处求助。
\end{enumerate}

\medskip\noindent\textbf{第三部分——讲解员的
技艺。}
\begin{enumerate}[resume]
  \item 参观者对模型发笑。你的开场白
        应像\cref{ch:g8:reproductive-systems}的
        第一段那样对待这两套系统。
        请写出这句话。
  \item 凭记忆，用
        \cref{met:g8:reproductive-systems:sides}中的
        五行对比概括两套系统。
  \item 有参观者问，婴儿的起点在道具上
        对应哪一部分。请将三件道具串起来，
        串成\cref{prop:g8:sexual-reproduction:core}
        所述的三个步骤。
  \item 用一句话结束开放日，说明
        健康中心为何还要讲授这部分解剖知识——
        以知识对抗传说，
        正如\cref{ex:g8:puberty:questions}所论证的。
\end{enumerate}
\end{problem}
